% Compiled with: latexmk -pdf -interaction=nonstopmode seminar.tex
\documentclass[aspectratio=169,11pt]{beamer}
\usetheme{default}
\usecolortheme{seagull}
\setbeamertemplate{navigation symbols}{}
\setbeamertemplate{footline}[frame number]
\usepackage{booktabs,graphicx,amsmath}

\title{Working Title}
\author{}
\institute{}
\date{Seminar draft}

\begin{document}

\begin{frame}\titlepage\end{frame}

\begin{frame}{Motivation}
% One question the audience already cares about; why now.
\end{frame}

\begin{frame}{This paper}
% Preview: what I do, the headline magnitude, the identification in one line.
\end{frame}

\begin{frame}{Data}
% Sources, sample window, unit of observation, one summary-stats table.
\end{frame}

\begin{frame}{The raw-data picture}
% ALWAYS before any regression table.
% \includegraphics[width=.85\textwidth]{../paper/figures/raw_fact.pdf}
\end{frame}

\begin{frame}{Empirical strategy}
% The estimating equation and the identifying assumption, nothing else.
\end{frame}

\begin{frame}{Baseline result}
% One result per slide. \input{../paper/tables/baseline.tex} (trim columns for slides)
\end{frame}

\begin{frame}{Threat 1 and its defense}
\end{frame}

\begin{frame}{Mechanism}
% Claims capped at "consistent with" absent direct evidence.
\end{frame}

\begin{frame}{Conclusion}
\end{frame}

\appendix
\begin{frame}{Appendix: robustness}
% One appendix slide per table a discussant might request.
\end{frame}

\end{document}
